\subsection{Checklist: Compute Resources}
\label{app:checklist_compute_resources}

\paragraph{Question.}
For each experiment, did we report the amount of compute and the type of
resources used?

\paragraph{Answer.}
Yes.

\paragraph{Justification.}
The paper-facing experiments consist of the \(15\)-system controlled
multibasin benchmark, the \(10\)-system Dysts \(dt{\times}30\) benchmark, and
several post-training evaluation diagnostics.  The main training jobs used the
cluster SLURM runners with CUDA enabled.  The scripts request one GPU for
training jobs, but do not pin a GPU model or constraint; the exact GPU type is
therefore cluster-assigned and should be treated as unspecified.  CPU-only
post-training analyses use SLURM CPU partitions or CPU-mode jobs.  The
estimates below are allocation upper bounds derived from the paper protocol
and the current SLURM resource requests, not measured wall-clock usage.

\paragraph{Hardware and resource details.}
Training jobs load CUDA \(12.6.0\) and request one cluster GPU, \(4\) CPU cores,
and \(16\) GB host memory per active training allocation.  Controlled
multibasin training uses single-run GPU allocations with a \(24\)-hour time
limit.  Dysts \(dt{\times}30\) training uses packed GPU allocations with up to
\(12\) runs per allocation, one GPU, \(4\) CPU cores, \(16\) GB memory, and a
\(72\)-hour time limit per packed allocation.  Dysts cache construction uses
CPU allocations with \(4\) CPU cores, \(16\) GB memory, and a \(24\)-hour time
limit.  Dysts long-horizon evaluation uses CPU allocations with \(4\) CPU
cores, \(24\) GB memory, and a \(6\)-hour time limit for the paper-facing
\(dt{\times}30\) evaluation queue.  Controlled support diagnostics and
support-routed analyses use CPU allocations with \(4\) CPU cores, \(16\)--\(24\)
GB memory, and \(8\)--\(12\)-hour time limits, depending on the diagnostic.

\paragraph{Per-run and total-compute estimates.}
\begin{center}
\begin{tabular}{|p{0.24\linewidth}|p{0.22\linewidth}|p{0.25\linewidth}|p{0.20\linewidth}|}
\hline
\textbf{Experiment} & \textbf{Paper-facing count} & \textbf{Per-run or per-allocation request} & \textbf{Allocation upper bound} \\
\hline
Controlled multibasin training & \(15\) systems \(\times\) \(6\) model rows \(\times\) \(15\) seeds \(=1350\) training runs & One GPU, \(4\) CPU cores, \(16\) GB memory, \(24\) h per run; \(200{,}000\) optimization steps, batch size \(256\), rollout length \(L=8\), latent dimension \(d_z=256\) & \(\le 32{,}400\) GPU-hours and \(\le 129{,}600\) allocated CPU-core-hours \\
\hline
Dysts \(dt{\times}30\) training & \(10\) systems \(\times\) \(6\) model rows \(\times\) \(15\) seeds \(=900\) training runs & Packed allocations: one GPU, \(4\) CPU cores, \(16\) GB memory, \(72\) h for up to \(12\) runs; \(100{,}000\) optimization steps, batch size \(256\), rollout length \(L=10\), latent dimension \(d_z=256\) & \(\le 5400\) GPU-hours and \(\le 21{,}600\) allocated CPU-core-hours when packs are full \\
\hline
Dysts cache construction and long-horizon evaluation & \(10\) retained systems; evaluation rows match the \(900\) trained Dysts runs & Cache jobs: \(4\) CPU cores, \(16\) GB, \(24\) h. Evaluation packs: \(4\) CPU cores, \(24\) GB, \(6\) h for up to \(12\) evaluation rows, plus one validation row and one collector & Approximately \(\le 4700\) CPU-core-hours for train/validation/test cache construction, packed evaluation, validation, and collection \\
\hline
Controlled support diagnostics for the main table & \(6\) model rows on the retained \(15\)-system benchmark & CPU reducer shards: \(4\) CPU cores, \(16\) GB, \(12\) h per model row; merge jobs: \(1\) CPU core, \(4\) GB, \(0.5\) h & Approximately \(\le 290\) CPU-core-hours for the six-row per-basin deep-slice support diagnostic pass \\
\hline
Appendix routed forecasting and support refresh diagnostics & Run on saved checkpoints; shard count depends on selected model rows and seed splits & CPU shards: \(4\) CPU cores, \(24\) GB, \(12\) h per model-row/seed-split shard; merges: \(1\) CPU core, \(4\) GB, \(0.5\) h & For \(6\) rows and three \(5\)-seed splits, \(\le 864\) CPU-core-hours per diagnostic family, plus merge overhead \\
\hline
Support-coordinate intervention figure & One representative trained checkpoint & One GPU, \(1\) CPU core, \(8\) GB memory, \(3\) h; plotting uses \(1\) CPU core, \(2\) GB, \(10\) min & \(\le 3\) GPU-hours, plus negligible CPU plotting time \\
\hline
\end{tabular}
\end{center}

\paragraph{Caveats.}
The totals above are conservative allocation ceilings.  Actual consumed compute
can be lower because SLURM jobs may finish before their time limits, completed
runs may be skipped, and packed jobs share one allocation across several
sequential runs.  Conversely, the repository contains older and supplemental
launchers for \(17\)-system controlled packets and \(12\)-system Dysts packets;
those extra runs are not included in the paper-facing totals above because the
manuscript reports the retained \(15\)-system controlled benchmark and the
retained \(10\)-system Dysts benchmark.  The scripts do not specify a GPU model,
so the hardware type is reported only as a single cluster-assigned CUDA GPU per
training allocation.
